% main_calculus_lecture2_metropolis.tex
\documentclass[aspectratio=169]{beamer}

% --- Language & encoding ---
\usepackage[utf8]{inputenc}
\usepackage[T1]{fontenc}
\usepackage[english,french]{babel}
\usepackage{lmodern}
\usepackage{microtype}

% --- Math & tables ---
\usepackage{amsmath,amssymb}
\usepackage{booktabs}
\usepackage{graphicx}
\usepackage{changepage}
\usepackage{amssymb}
\newcommand{\R}{\mathbb{R}}
% --- Metropolis Theme ---
\usetheme[numbering=fraction,progressbar=frametitle]{metropolis}
\setbeamertemplate{frame numbering}[fraction]

% --- Colors (keep your palette) ---
\definecolor{BootcampBlueDark}{HTML}{1A3C68}
\definecolor{TextBlack}{HTML}{000000}
\definecolor{AccentGold}{HTML}{DAA520}
\definecolor{Crimson}{HTML}{990000}
\definecolor{MatrixGreen}{HTML}{228B22}
\definecolor{MatrixRed}{HTML}{B22222}

\setbeamercolor{block title example}{bg=BootcampBlueDark!90!black,fg=white}
\setbeamercolor{block body example}{bg=BootcampBlueDark!5!white,fg=black}

% === Thick frame title band ===
\setbeamercolor{frametitle}{bg=BootcampBlueDark,fg=white}
\setbeamerfont{frametitle}{series=\bfseries,size=\large}
\setbeamertemplate{frametitle}{
  \nointerlineskip
  \begin{beamercolorbox}[wd=\paperwidth,ht=3.2ex,dp=1.4ex,leftskip=1em,rightskip=1em]{frametitle}
    \usebeamerfont{frametitle}\insertframetitle
  \end{beamercolorbox}
}

% === Custom Definition Block ===
\newenvironment{definitionblock}[1]{%
  \par\vspace{0.4em}%
  \begin{beamercolorbox}[
      wd=\textwidth,
      sep=0.4em,
      leftskip=0.3em,
      rightskip=0.6em,
      rounded=false,
      shadow=false
    ]{block title example}%
    \raisebox{0.15em}{\usebeamerfont{block title example}#1}%
  \end{beamercolorbox}%
  \vspace{-0.4em}%
  \begin{beamercolorbox}[
      wd=\textwidth,
      sep=1em,
      leftskip=0.6em,
      rightskip=0.6em,
      rounded=false,
      shadow=false
    ]{block body example}%
    \usebeamerfont{block body example}%
}{%
  \end{beamercolorbox}%
  \par\vspace{0.6em}%
}

% === Global text formatting ===
\setbeamercolor{normal text}{fg=TextBlack,bg=white}
\setbeamercolor{title}{fg=TextBlack}
\setbeamercolor{structure}{fg=BootcampBlueDark}
\setbeamercolor{progress bar}{fg=BootcampBlueDark,bg=gray!30}

% --- Course Info ---
\title{Lecture 2 — Differential Calculus (1D $\to$ Multivariate), Integration, Taylor, \& Optimization Links}
\institute{NYU Tandon School, Finance and Risk Engineering}

\metroset{
  sectionpage=progressbar,
  subsectionpage=progressbar
}

% --- Section & subsection transitions (plain) ---
\AtBeginSection{\frame[plain]{\sectionpage}}
\AtBeginSubsection{\frame[plain]{\subsectionpage}}

\begin{document}

\maketitle

\begin{frame}{Outline}
  \tableofcontents[hideallsubsections]
\end{frame}

% ==================================================
% SECTION 1 — Differential Calculus (1D → nD)
% ==================================================
\section{Differential Calculus: From One to Several Variables}

\begin{frame}{Outline}
  \tableofcontents[currentsection,hideothersubsections]
\end{frame}

% ======================================================
\subsection{Univariate Derivatives \& Rules}
% ======================================================

\begin{frame}{Derivative: Definition, Interpretation, Notation}
\begin{definitionblock}{Derivative (limit definition)}
\[
f'(\textcolor{AccentGold}{x})
=\lim_{h\to 0}
\dfrac{f(\textcolor{AccentGold}{x}+h)-f(\textcolor{AccentGold}{x})}{h},
\qquad\text{if the limit exists}.
\]
\end{definitionblock}

\begin{itemize}
    \item Measures the \textbf{instantaneous rate of change} of \(f\) at \(x\).
    \item \textbf{Slope} of the tangent line at the point \((x,f(x))\).
    \item Notations used:
    \[
    f'(x),\quad \dfrac{df}{dx}(x),\quad \partial_x f(x).
    \]
    \item Differentiability \(\Rightarrow\) continuity, but not conversely.
\end{itemize}
\end{frame}

\begin{frame}{Differentiation Rules (Complete List)}
\begin{itemize}
    \item \textbf{Linearity}
    \[
    (af+bg)' = af' + bg', \qquad a,b\in\mathbb{R}.
    \]

    \item \textbf{Product rule}
    \[
    (fg)' = f'g + fg'.
    \]
    Example: \( (x^2 e^x)' = 2x e^x + x^2 e^x\).

    \item \textbf{Quotient rule}
    \[
    \left(\frac{f}{g}\right)' = 
    \frac{f'g - fg'}{g^2}, \qquad g(x)\neq 0.
    \]
    Example: \( \left(\frac{x}{x+1}\right)' = \frac{1}{(x+1)^2}. \)

    \item \textbf{Chain rule}
    If \(h(x)=f(g(x))\), then
    \[
    h'(x)=f'(g(x))\,g'(x).
    \]
    Example: \( \dfrac{d}{dx}\sin(x^2)=\cos(x^2)\cdot 2x.\)
\end{itemize}
\end{frame}

\begin{frame}{Classic Derivatives (Useful Table)}
\begin{itemize}
    \item Powers: \( \dfrac{d}{dx}(x^n)=n x^{n-1}\).
    \item Exponential: \( (e^x)' = e^x \).
    \item Logarithm: \( (\ln x)' = \dfrac{1}{x},\quad x>0.\)
    \item Trigonometric:
    \[
    (\sin x)'=\cos x,\qquad (\cos x)'=-\sin x.
    \]
    \item Inverse functions : if $ f \circ f^{-1} (x) = x$, then $ \left. f^{-1} \right.' (x) = \frac{1}{ f' \circ f^{-1}(x)}.$
\end{itemize}
\end{frame}

\begin{frame}{Worked Example: Chain Rule with Polynomial Composition}
\textbf{Goal:} compute the derivative of
\[
F(x)=(2x+1)^4.
\]

\begin{itemize}
    \item Let \(u(x)=2x+1\).
    \item Then \(F(x)=u(x)^4\).
\end{itemize}

\[
F'(x)=4\,u(x)^3\cdot u'(x)
     =4(2x+1)^3\cdot 2.
\]

\[
\boxed{F'(x)=8(2x+1)^3}.
\]
\end{frame}

% ======================================================
\subsection{Partial Derivatives \& Gradient}
% ======================================================

\begin{frame}{From Partial Derivatives to the Gradient (Multivariate)}

\begin{definitionblock}{Partial derivatives}
For \(f:\mathbb{R}^n\to\mathbb{R}\), the partial derivative with respect to \(x_i\) is
\[
\partial_{x_i} f(x)
=
\lim_{h\to 0}
\frac{f(x_1,\dots,x_i+h,\dots,x_n)-f(x)}{h},
\qquad i=1,\dots,n.
\]
We then define the \textbf{Gradient} as
$\nabla f(x)
=
\begin{pmatrix}
\partial_{x_1}f(x)\\[-2pt]
\vdots\\[-2pt]
\partial_{x_n}f(x)
\end{pmatrix}
\in\mathbb{R}^n.
$
We then have $\nabla f : \R^n \rightarrow \R^n$.
\end{definitionblock}
\vspace{-5mm}
\begin{itemize}
    \item \textbf{Geometric meaning:}  
    \(\nabla f(x)\) points in the direction of **steepest ascent**;  
    \(\|\nabla f(x)\|\) is the maximum directional rate of change at \(x\).
    \item \textbf{Directional derivative:} $
 D_{\!u}f(x)=\nabla f(x)^\top u,\qquad \|u\|=1.$
\end{itemize}

\end{frame}

\begin{frame}{Example: Gradient of a Quadratic Form}

Let
\[
f(x)=x^\top Q x,
\qquad Q\in\mathbb{R}^{n\times n},\; x\in\mathbb{R}^n.
\]

\textbf{Goal: compute the gradient.}

\medskip

Expand: $
f(x)=\sum_{i=1}^n\sum_{j=1}^n Q_{ij} x_i x_j.
$
Differentiate with respect to \(x_k\):
\[
\partial_{x_k} f(x)
=
\sum_{j=1}^n Q_{kj} x_j
+
\sum_{i=1}^n Q_{ik} x_i.
\]
Thus, in vector form,
\[
\nabla f(x) = (Q+Q^\top)x.
\]

\begin{itemize}
    \item If \(Q\) is symmetric, \(\nabla f(x)=2Qx\).
    \item This example will be essential when we study \textbf{convexity} and \text{Hessian}.
\end{itemize}

\end{frame}

\subsection{Vector-Valued Functions \& Jacobian}

\begin{frame}{Vector-Valued Functions \& the Jacobian Matrix}

\begin{definitionblock}{Vector-Valued Function}
A function \(g:\mathbb{R}^n \to \mathbb{R}^m\) maps an \(n\)-dimensional input
vector to an \(m\)-dimensional output vector:
\[
g(x)=
\begin{pmatrix}
g_1(x)\\
\vdots\\
g_m(x)
\end{pmatrix},
\qquad
x=
\begin{pmatrix}
x_1\\
\vdots\\
x_n
\end{pmatrix}.
\]
\end{definitionblock}

\begin{itemize}
    \item Each component \(g_i(x)\) may depend on all variables \(x_1,\ldots,x_n\).
    \item Examples: $
        g(x_1,x_2)=\begin{pmatrix}x_1^2+x_2\\ e^{x_1-x_2}\end{pmatrix}
        \,\, (n=2,\ m=2)$ and $
        h(x_1,x_2,x_3) =\begin{pmatrix}x_1+x_2+x_3 \\ x_1x_2\end{pmatrix}
        \,\, (n=3,\ m=2).$
\end{itemize}

\end{frame}

\begin{frame}{Jacobian Matrix \(J_g(x)\)}

\begin{definitionblock}{Jacobian of \(g:\mathbb{R}^n\to\mathbb{R}^m\)}
\[
J_g(x)=\frac{\partial g(x)}{\partial x}=
\begin{bmatrix}
\dfrac{\partial g_1}{\partial x_1} & \cdots & \dfrac{\partial g_1}{\partial x_n}\\
\vdots & \ddots & \vdots\\
\dfrac{\partial g_m}{\partial x_1} & \cdots & \dfrac{\partial g_m}{\partial x_n}
\end{bmatrix}
\in\mathbb{R}^{m\times n}.
\]
\end{definitionblock}

\begin{itemize}
    \item \(J_g(x)\) is the **matrix of all first-order partial derivatives**.
    \item Linear approximation around a point \(x_0\):
    \[
    g(x_0+h) \approx g(x_0) + J_g(x_0)h.
    \]
    So \(J_g(x_0)\) generalizes the slope to multiple dimensions.
\end{itemize}

\end{frame}

\begin{frame}{Geometric \& Analytic Interpretation}

\begin{itemize}
    \item \(J_g(x)\) acts as the **best linear map** approximating \(g\) near \(x\):
    \[
    g(x+h) \approx g(x) + J_g(x)\,h.
    \]

    \item The Jacobian describes:
    \begin{itemize}
        \item How small changes in inputs affect outputs.
        \item How a surface/volume in \(\mathbb{R}^n\) is stretched or rotated under \(g\).
    \end{itemize}

    \item For \(m=n\), \(|\det J_g(x)|\) gives the **local volume-scaling factor**:
    \[
    \text{Vol}(g(A)) \approx |\det J_g(x)|\,\text{Vol}(A)
    \]
\end{itemize}

\end{frame}

\begin{frame}{Chain Rule (Matrix Form) \& Example}

\textbf{Chain rule (matrix version).}
If
\[
f:\mathbb{R}^m \to \mathbb{R}^p, \qquad
g:\mathbb{R}^n \to \mathbb{R}^m,
\]
then their composition \(h=f\circ g:\mathbb{R}^n\to\mathbb{R}^p\) satisfies
\[
J_h(x)
= J_f(g(x))\, J_g(x).
\]
\textit{(Matrix multiplication, dimension check: \(p\times m\) by \(m\times n\) yields \(p\times n\)).}

\medskip

\textbf{Example.}
Let \(g(x_1,x_2)=(x_1^2+x_2,\, e^{x_1-x_2})\) and \(f(y_1,y_2)=y_1+y_2^2\). We have 
$
J_g(x)
=
\begin{bmatrix}
2x_1 & 1\\
e^{x_1-x_2} & -e^{x_1-x_2}
\end{bmatrix}$, $
J_f(y)
=
\begin{bmatrix}
1 & 2y_2
\end{bmatrix}$.
Then
$
J_h(x)
=J_f(g(x))J_g(x)$, then 
$
J_h(x)
=\begin{bmatrix}
1 & 2e^{x_1-x_2}
\end{bmatrix}
\begin{bmatrix}
2x_1 & 1\\
e^{x_1-x_2} & -e^{x_1-x_2}
\end{bmatrix}
$ and 
$
\boxed{J_h(x)
=\begin{bmatrix}
2x_1+2e^{2(x_1-x_2)} &
1-2e^{2(x_1-x_2)}
\end{bmatrix}.}
$

\end{frame}

\begin{frame}{Examples \& Applications}

\textbf{Linear map:}  
If \(g(x)=Ax+b\), \(A\in\mathbb{R}^{m\times n}\), then
\[
J_g(x)=A.
\]
The Jacobian equals the transformation matrix itself.

\medskip

\textbf{Nonlinear transformation (polar coordinates).}
\[
g(r,\theta)=
\begin{pmatrix}
r\cos\theta\\[2pt]
r\sin\theta
\end{pmatrix}.
\qquad
J_g(r,\theta)=
\begin{bmatrix}
\cos\theta & -r\sin\theta\\
\sin\theta & r\cos\theta
\end{bmatrix}.
\]
\[
\det J_g(r,\theta)=r,
\]
the classical **Jacobian determinant** in change of variables
\((x,y)\leftrightarrow(r,\theta)\).
\end{frame}

% ======================================================
\subsection{Second Derivatives \& Schwarz Theorem}
% ======================================================

\begin{frame}{Second Mixed Derivatives \& Regularity}

\begin{definitionblock}{Second Mixed Partial Derivatives}
For a function \( f:\mathbb{R}^n \to \mathbb{R} \),
\[
\partial_{x_i}\partial_{x_j} f(x)
=
\frac{\partial}{\partial x_i}
\!\left(\frac{\partial f}{\partial x_j}(x)\right),
\qquad i,j=1,\dots,n.
\]
\end{definitionblock}


\medskip
\textbf{Notation:}\quad
\( f_{x_i x_j}(x) = \partial_{x_i}\partial_{x_j} f(x). \)

\end{frame}

\begin{frame}{Schwarz Theorem}

\begin{definitionblock}{Schwarz (Symmetry of Mixed Partials)}
If \( f \) has continuous second-order partial derivatives in a neighborhood of \(x\), then
\[
\boxed{
\partial_{x_i}\partial_{x_j} f(x)
=
\partial_{x_j}\partial_{x_i} f(x),
\qquad \forall i,j.
}
\]
\end{definitionblock}

\begin{itemize}
    \item Key assumption: \( f \in C^2 \) (all second derivatives exist and are continuous).
\end{itemize}

\medskip
\textbf{Practical implication:}
Guarantees the consistency of quadratic (second-order Taylor) approximations used in optimization and PDE models.

\end{frame}

\begin{frame}{Hessian Matrix: Definition \& Structure}

\begin{definitionblock}{Hessian of \( f:\mathbb{R}^n \to \mathbb{R} \)}
\[
\nabla^2 f(x)=
\begin{pmatrix}
\partial_{x_1x_1}f(x) & \cdots & \partial_{x_1x_n}f(x)\\
\vdots & \ddots & \vdots\\
\partial_{x_nx_1}f(x) & \cdots & \partial_{x_nx_n}f(x)
\end{pmatrix}
\in \mathbb{R}^{n\times n}.
\]
\end{definitionblock}

\begin{itemize}
    \item If \( f \in C^2 \), then \( \nabla^2 f(x) \) is \textbf{symmetric} (by Schwarz Theorem).
    \item The Hessian captures the **local curvature** of \(f\).
    \item Central tool for:
    \begin{itemize}
        \item Optimization (minima, maxima, saddle points),
        \item Convexity analysis,
        \item Second-order Taylor approximations,
        \item Quadratic PDEs and control problems.
    \end{itemize}
\end{itemize}

\end{frame}

% ----------------------------------------------------------
\begin{frame}{Example: Hessian of a Quadratic Form}

Consider the quadratic function
\[
f(x) = x^\top Q x, \qquad x \in \mathbb{R}^n, \; Q \in \mathbb{R}^{n\times n}.
\]

\begin{itemize}
    \item Gradient:
    \[
    \nabla f(x) = (Q + Q^\top)x.
    \]
    \item Hessian:
    \[
    \nabla^2 f(x) = Q + Q^\top.
    \]
\end{itemize}

\textbf{Special case:} if \(Q\) is symmetric,
\[
\nabla f(x) = 2Qx, 
\qquad
\nabla^2 f(x) = 2Q.
\]

\begin{itemize}
    \item The Hessian is constant (does not depend on \(x\)).
    \item Convexity test: \(f\) is convex iff \(Q\succeq 0\).
\end{itemize}

\end{frame}

% ==================================================
% SECTION 2 — Integration
% ==================================================
\section{Integration}

\begin{frame}{Outline}
  \tableofcontents[currentsection,hideothersubsections]
\end{frame}

% ============================================
\subsection{Recall: 1D Integrals}
% ============================================

\begin{frame}{1D Integrals: Definition \& Interpretation}

\begin{definitionblock}{Definite Integral}
For an integrable function \(f:[a,b]\to\mathbb{R}\), the definite integral
$
\int_a^b f(x)\,dx
$
is the signed area under the curve.
\end{definitionblock}

\begin{itemize}
    \item Measures cumulative quantity (area, mass, probability).
    \item Negative contributions when \(f(x)<0\).
    \item Additivity:
    \[
    \int_a^c f = \int_a^b f + \int_b^c f.
    \]
    \item Linearity:
    \[
    \int_a^b (\alpha f + \beta g)
    = \alpha \int_a^b f + \beta \int_a^b g.
    \]
\end{itemize}

\end{frame}

\begin{frame}{Fundamental Theorem of Calculus (FTC)}

\begin{definitionblock}{Fundamental Theorem of Calculus}
If \(f\) is continuous on \([a,b]\) and
\[
F(x)=\int_a^x f(t)\,dt,
\]
then
\[
F'(x)=f(x), \qquad \text{and} \qquad
\int_a^b f(x)\,dx = F(b)-F(a).
\]
\end{definitionblock}

\begin{itemize}
    \item Antiderivative \(F\) recovers the integral.
    \item Differentiation and integration are inverse operations (under regularity).
\end{itemize}

\end{frame}

\begin{frame}{Change of Variables (1D)}

\begin{definitionblock}{Change of Variables Formula}
Let \(u=g(x)\) with \(g\) strictly monotone and differentiable. Then
\[
\int_{a}^{b} f(x)\,dx
=
\int_{g(a)}^{g(b)} f(g^{-1}(u))\,
\left|(g^{-1})'(u)\right|\,du.
\]
Equivalent formulation:
\vspace{-3mm}
\[
\int f(g(x))\,g'(x)\,dx = \int f(u)\,du.
\vspace{-6mm}
\]
\end{definitionblock}
\vspace{-6 mm}
\begin{itemize}
    \item Most useful when one substitution simplifies the integrand.
    \item Typical example: Gaussian integrals, logarithmic/exponential substitutions.
\end{itemize}

\textbf{Example (simple):}  
Let \(u = 1+x^2\), then \(du = 2x\,dx\):
\[
\int_0^1 \frac{2x}{1+x^2}\,dx
=
\int_1^2 \frac{1}{u}\,du
= \ln 2.
\]

\end{frame}

\begin{frame}{Integration by Parts}

\begin{definitionblock}{Integration by Parts}
If \(u=u(x)\) and \(v=v(x)\) are differentiable:
\vspace{-4mm}
\[
\int_a^b u(x) v'(x)\,dx
=
\big[u(x)v(x)\big]_{a}^{b}
-
\int_a^b u'(x) v(x)\,dx.
\vspace{-6mm}
\]
\end{definitionblock}
\vspace{-3mm}
\begin{itemize}
    \item Differentiation “moves” between factors.
    \item Often used with:
    \begin{itemize}
        \item polynomials × exponentials,
        \item polynomials × trig functions,
        \item logarithms.
    \end{itemize}
\end{itemize}

\textbf{Example:}
\[
\int_0^1 x e^x\,dx
=
\big[x e^x\big]_0^1 - \int_0^1 1\cdot e^x\,dx
=
e - (e-1) = 1.
\]

\end{frame}

\begin{frame}{Worked Example: Expectation Under a Simple Density}

Let \(X\) have density
\[
f_X(x) = 2x, \qquad x\in[0,1].
\]

\textbf{Check normalization:}
\[
\int_0^1 2x\,dx = 1.
\]

\begin{definitionblock}{Expectation}
\[
\mathbb{E}[X] = \int_0^1 x f_X(x)\,dx = \int_0^1 2x^2\,dx.
\]
\end{definitionblock}

\[
\mathbb{E}[X]
= 2\left[\frac{x^3}{3}\right]_{0}^{1}
= \frac{2}{3}.
\]

\begin{itemize}
    \item Integration computes probability, expectation, and variance.
    \item Same rules apply for any 1D continuous random variable.
\end{itemize}

\end{frame}

% ======================================================
\subsection{Leibniz Rule (Differentiate Under the Integral)}
% ======================================================

\begin{frame}{Leibniz Rule: General Case (Parameter \(\theta\))}

\begin{definitionblock}{Leibniz Rule (general form)}

Let \(f(x,\theta)\) be a function defined on a domain  
$
\{(x,\theta): a(\theta)\le x\le b(\theta),\ \theta\in I\},
$
and assume:

\begin{enumerate}
    \item \(a(\theta), b(\theta)\) are differentiable on \(I\);
    \item \(f(\cdot,\theta)\) is integrable for each \(\theta\);
    \item \(f(x,\theta)\) and \(\partial_\theta f(x,\theta)\) are continuous on the domain;
    \item There exists an integrable dominating function \(M(x)\) such that  
    \(|\partial_\theta f| \le M(\) for all \(\theta\in I\)
    \ (Domination ensures one may interchange differentiation and integration).
\end{enumerate}

Then:
\[
\frac{d}{d\theta}\int_{a(\theta)}^{b(\theta)} f(x,\theta)\,dx
=
f(b(\theta),\theta)\,b'(\theta)
-
f(a(\theta),\theta)\,a'(\theta)
+
\int_{a(\theta)}^{b(\theta)}\frac{\partial f}{\partial \theta}(x,\theta)\,dx.
\]

\end{definitionblock}

\begin{itemize}
    \item This is a consequence of the Dominated Convergence Theorem.  
    \item Needed whenever both the integrand and the limits depend on a parameter.
    \item Core tool for sensitivity analysis (Greeks), parametric expectations, and statistical models.
\end{itemize}

\end{frame}

\begin{frame}{Special Case: Fixed Limits}

If \(a(\theta)=a\) and \(b(\theta)=b\) are constants, then
\[
\frac{d}{d\theta} \int_a^b f(x,\theta)\,dx
=
\int_a^b \frac{\partial f}{\partial \theta}(x,\theta)\,dx.
\]

\begin{itemize}
    \item The derivative “passes inside’’ the integral.
    \item Holds under mild regularity (continuity of the mixed partial).
    \item Used in machine learning, statistics, Monte Carlo, physics (Feynman).
\end{itemize}

\medskip

\textbf{Example (simple):}
\[
I(\theta)=\int_0^1 e^{-\theta x}\,dx.
\]

Then
\[
\frac{d}{d\theta}I(\theta)
=
\int_0^1 \frac{\partial}{\partial \theta}e^{-\theta x}\,dx
=
\int_0^1 -x e^{-\theta x}\,dx.
\]

\end{frame}

\begin{frame}{The “Feynman Trick”}

\textbf{Idea:}  
Introduce a parameter \(\theta\) to make the integral easier to differentiate,  
then integrate the result.

\medskip
\textbf{Classical example:}
$I = \int_{\R^+} \frac{\sin(x)}{x} \mathrm{d}x$ can be computed by considering $I(a) =  \int_{\R^+} e^{-ax}\frac{\sin(x)}{x} \mathrm{d}x$
\begin{itemize}
    \item Famous technique in probability.
\end{itemize}

\end{frame}

\begin{frame}{Leibniz Rule with Moving Boundaries: Worked Example}

Let
$
I(\theta)=\int_0^{\theta} \sin(\theta x)\,dx.
$
Here:
$
a(\theta)=0$, $b(\theta)=\theta$ and, $f(x,\theta)=\sin(\theta x).
$
Apply Leibniz:
\[
\frac{dI}{d\theta}
=
f(\theta,\theta)\cdot b'(\theta)
-
f(0,\theta)\cdot a'(\theta)
+
\int_{0}^{\theta}
\frac{\partial}{\partial \theta}\sin(\theta x)\,dx.
\]
Compute:
\[
f(\theta,\theta)=\sin(\theta^2), \qquad
f(0,\theta)=\sin(0)=0,
\]
\[
\frac{\partial}{\partial \theta}\sin(\theta x)=x\cos(\theta x).
\]

Thus
\[
\frac{dI}{d\theta}
=
\sin(\theta^2)
+
\int_{0}^{\theta} x\cos(\theta x)\,dx.
\]

(Integral can be evaluated by substitution \(u=\theta x\), or left symbolic.)

\end{frame}

\begin{frame}{Application: Sensitivity of Expectations}

If
\[
F(\theta) = \mathbb{E}_\theta[g(X,\theta)]
=
\int_{\mathbb{R}} g(x,\theta)\,f_\theta(x)\,dx,
\]
then under mild assumptions,
\[
\frac{dF}{d\theta}
=
\int_{\mathbb{R}}
\left(
\frac{\partial g}{\partial \theta}(x,\theta)
+
g(x,\theta)
\frac{\partial}{\partial \theta}\log f_\theta(x)
\right) f_\theta(x)\,dx.
\]

\begin{itemize}
    \item Used to compute Greeks in quantitative finance.
    \item Key to likelihood derivatives (score function) in statistics.
    \item Foundation of the “log-derivative trick’’ in reinforcement learning.
\end{itemize}

\medskip

\textbf{Example (Greek-like sensitivity):}
\[
V(\sigma)=\int_0^\infty \text{Payoff}(x)\,p(x;\sigma)\,dx,
\qquad
\frac{dV}{d\sigma}
=
\int_0^\infty \text{Payoff}(x)\,\frac{\partial p}{\partial \sigma}(x;\sigma)\,dx.
\]

\end{frame}

% ======================================================
\subsection{n-Dimensional Integration \& Vector Calculus Identities}
% ======================================================

\begin{frame}{Integration in \(\mathbb{R}^n\): Basic Definition}

\begin{definitionblock}{Integral of \(f:\mathbb{R}^n\to\mathbb{R}\)}
Given a measurable region \(D\subset\mathbb{R}^n\) and an integrable function \(f\),
\[
\int_D f(x)\,dx
=
\lim_{\text{mesh}\to 0}
\sum_{k} f(\xi_k)\,\mathrm{Vol}(\Delta_k).
\]
\end{definitionblock}

\begin{itemize}
  \item Generalizes Riemann sums: volume elements \(\mathrm{Vol}(\Delta_k)\) replace lengths/areas.
  \item In practice: iterated integrals via Fubini (when integrable).
  \item Volume element: \(dx = dx_1\cdots dx_n\).
  \item Often evaluated after a suitable change of variables.
\end{itemize}
\end{frame}

\begin{frame}{Fubini: Iterated Integrals in \(\mathbb{R}^n\)}

\begin{definitionblock}{Fubini's Theorem (informal)}
If \(f\) is integrable on a rectangular box \(D\),
\[
\int_D f(x_1,\dots,x_n)\,dx
=
\int_{a_1}^{b_1}\!\!\cdots\!\int_{a_n}^{b_n}
f(x_1,\dots,x_n)\,dx_n\cdots dx_1.
\]
\end{definitionblock}

\begin{itemize}
  \item Practical computational tool.
  \item Important for probability: joint densities, expectations.
  \item Valid for more general “nice’’ regions with indicator functions.
\end{itemize}
\end{frame}

\begin{frame}{Change of Variables in \(\mathbb{R}^n\)}

\begin{definitionblock}{Change of Variables}
Let \(x = \Phi(u)\) be a \(C^1\) bijection with invertible Jacobian. Then
\[
\int_{\Phi(U)} f(x)\,dx
= 
\int_{U} f(\Phi(u))\,\big|\det J_\Phi(u)\big|\,du.
\]
\end{definitionblock}

\begin{itemize}
  \item The Jacobian determinant measures local volume distortion.
  \item Use to simplify complicated regions or integrands.
  \item Spherical, polar, cylindrical coordinates = special cases.
\end{itemize}
\end{frame}

\begin{frame}{Jacobian Determinant: Geometry}

\begin{itemize}
  \item \(J_\Phi(u)\): matrix of first-order partials.
  \item \(|\det J_\Phi(u)|\): volume of the parallelepiped spanned by transformed basis vectors.
\end{itemize}

\[
|\det J_\Phi(u)| = \frac{\text{volume of small box in }x\text{-space}}
{\text{volume of small box in }u\text{-space}}.
\]

\textbf{Examples}
\begin{itemize}
  \item Polar: \(|\det J| = r\).
  \item Spherical: \(|\det J| = r^2 \sin\theta\).
\end{itemize}
\end{frame}

\begin{frame}{Example: Gaussian Integral via Change of Variables}

Goal:
$
I = \int_{0}^{\infty} e^{-x^{2}}\,dx.
$
Square the integral:
\[
I^2 = \left(\int_0^\infty e^{-x^2}\,dx\right)
      \left(\int_0^\infty e^{-y^2}\,dy\right)
    = \int_{0}^{\infty}\!\!\int_{0}^{\infty} e^{-(x^2+y^2)}\,dx\,dy.
\]

\textbf{Switch to polar coordinates: } 
 Letting $J(r, \theta) = ( r \cos (\theta), r \sin(\theta))$, we have for \(x = r\cos\theta,\ y = r\sin\theta\), $\det J = r.$
Thus
\[
I^2 = \int_{0}^{\pi/2}\!\int_{0}^{\infty}
      e^{-r^2}\,r\,dr\,d\theta
    = \left(\int_0^{\pi/2} d\theta\right)
      \left(\int_{0}^{\infty} r e^{-r^2}\,dr\right).
\]

Compute:
$
\int_{0}^{\pi/2} d\theta = \frac{\pi}{2},
$, $
\int_0^\infty r e^{-r^2}\,dr = \frac12.
$
Therefore
\[
I^2 = \frac{\pi}{2}\cdot\frac12 = \frac{\pi}{4},
\qquad 
I = \frac{\sqrt{\pi}}{2}.
\]

\textbf{Key principle:} multivariate change of variables simplifies otherwise intractable 1D integrals.
\end{frame}


% ==================================================
% SECTION 3 — Taylor Expansion + Classical Inequalities
% ==================================================
\section{Taylor Expansion \& Classical Inequalities}

\begin{frame}{Outline}
  \tableofcontents[currentsection,hideothersubsections]
\end{frame}

% ======================================================
\subsection{Taylor (1D): Polynomial, Series, Remainder}
% ======================================================

\begin{frame}{Taylor Expansion in One Dimension}

\begin{definitionblock}{Taylor Polynomial of Degree \(n\) at \(x_0\)}
If \(f\) has \(n\) derivatives around \(x_0\), its Taylor polynomial is
\[
T_n(x)
=
\sum_{k=0}^n 
\frac{f^{(k)}(x_0)}{k!}\,(x-x_0)^k.
\]
\end{definitionblock}

\begin{itemize}
  \item Local approximation of \(f\) near \(x_0\).
  \item Linear approx = 1st-order; quadratic approx = 2nd-order.
  \item Used heavily in optimization, risk sensitivities, PDEs.
\end{itemize}

\end{frame}

\begin{frame}{Taylor Series \& Maclaurin Expansion}

\textbf{Infinite-order Taylor series:}
\[
f(x) = \sum_{k=0}^{\infty} \frac{f^{(k)}(x_0)}{k!}(x-x_0)^k
\quad\text{(if convergent)}.
\]

\begin{itemize}
  \item Not all series converge; not all convergent series equal \(f\).
  \item Important notions:
        \begin{itemize}
            \item Radius of convergence \(R\) (from power series theory).
            \item Interval where polynomial approximations are accurate.
        \end{itemize}
  \item Canonical examples: \(\exp(x)\), \(\log(1+x)\), \(\sin x\), \(\cos x\).
\end{itemize}

% TODO: add a simple expansion: e.g., $\exp(x)=1+x+x^2/2+\cdots$.

\end{frame}

\begin{frame}{Integral Form of the Remainder}

\begin{definitionblock}{Integral Remainder}
If \(f^{(n+1)}\) is continuous on an interval containing \(x_0,x\), then
\[
R_{n+1}(x)
=
f(x) - T_n(x)= 
\frac{1}{n!}\int_{x_0}^{x} (x-t)^n f^{(n+1)}(t)\,dt.
\]
\end{definitionblock}

\begin{itemize}
  \item Equivalent to the Lagrange form by the mean value theorem for integrals.
  \item \textbf{Useful for analysis:} allows estimates via integral norms.
  \item If \(|f^{(n+1)}(t)|\le M\), we immediately obtain the bound
  \[
  |R_{n+1}(x)|\le \frac{M}{(n+1)!}|x-x_0|^{n+1}.
  \]
\end{itemize}

\end{frame}

% ------------------------------------------------------

\begin{frame}{Geometric \& Analytical Interpretation}

\begin{itemize}
  \item The remainder measures how much the function’s true value
        deviates from its Taylor polynomial.
  \item \textbf{Graphically:} it is the “vertical gap’’ between the curve
        and its local polynomial approximation.
  \item The order of the remainder tells how quickly this gap shrinks:
        \[
        R_{n+1}(x) = o((x-x_0)^n)
        \quad\Rightarrow\quad
        \frac{R_{n+1}(x)}{(x-x_0)^n}\to 0.
        \]
  \item Smooth functions (\(C^\infty\)) have remainders of arbitrarily high order,
        forming analytic series when convergence holds.
\end{itemize}

\end{frame}

% ------------------------------------------------------

\begin{frame}{Examples \& Applications}

\textbf{Example 1 — Exponential function:}
\[
e^x = 1 + x + \frac{x^2}{2!} + \cdots + \frac{x^n}{n!}
      + R_{n+1}(x),
\qquad
R_{n+1}(x)=\frac{e^\xi}{(n+1)!}x^{n+1}.
\]
Since \(e^\xi \le e^{|x|}\),
\[
|R_{n+1}(x)| \le \frac{e^{|x|}}{(n+1)!}|x|^{n+1}.
\]

\textbf{Example 2 — Sine function:}
\[
\sin x = x - \frac{x^3}{3!} + R_5(x),
\qquad
|R_5(x)| \le \frac{|x|^5}{5!}.
\]

\textbf{Application:}
Error control for truncated series, numerical methods, sensitivity estimates in Greeks, or confidence-interval expansions.

\end{frame}

% ------------------------------------------------------

\begin{frame}{Summary: Remainder Forms Compared}

\begin{table}[]
\centering
\begin{tabular}{l|l}
\textbf{Form} & \textbf{Expression for } \(R_{n+1}(x)\) \\
\hline
\textbf{Lagrange} &
\(\displaystyle \frac{f^{(n+1)}(\xi)}{(n+1)!}(x-x_0)^{n+1}\) \\[4pt]
\textbf{Integral} &
\(\displaystyle \frac{1}{n!}\int_{x_0}^{x}(x-t)^n f^{(n+1)}(t)\,dt\) \\[6pt]
\textbf{Peano} &
\(o((x-x_0)^n)\) as \(x\to x_0\) \\[4pt]
\end{tabular}
\end{table}

\begin{itemize}
  \item All forms describe the same concept at increasing levels of generality.
  \item The Lagrange form is most useful for explicit numerical bounds.
  \item The integral and Peano forms are suited for theoretical analysis.
\end{itemize}

\end{frame}

% ======================================================
\subsection{Multivariate Taylor (to 2nd order)}
% ======================================================

\begin{frame}{Multivariate Taylor \& Linearization}

\begin{definitionblock}{Taylor Expansion of \(f:\mathbb{R}^n\to\mathbb{R}\) at \(x\)}
For small \(h\in\mathbb{R}^n\),
\[
f(x+h)
=
f(x)
+ \nabla f(x)^\top h
+ \frac12\, h^\top \nabla^2 f(x)\, h
+ R_3(h),
\]
where \(R_3(h) = o(\|h\|^2)\) as \(\|h\|\to 0\).
\end{definitionblock}
\vspace{-5mm}
\begin{itemize}
  \item \textbf{1st-order approximation:}
    $
    f(x+h) \approx f(x) + \nabla f(x)^\top h.
    $
    Linearization / sensitivity of \(f\) to small perturbations.
  \item \textbf{2nd-order (quadratic) approximation:}
    \[
    f(x+h) \approx f(x) + \nabla f(x)^\top h
    + \tfrac12 h^\top H(x) h.
    \]
    Incorporates curvature information through the Hessian.
  \item Crucial in optimization (Newton methods, convexity tests).
\end{itemize}

% TODO: add example with quadratic form or log-sum-exp.

\end{frame}

\begin{frame}{Interpretation of the Terms}

\textbf{Zeroth order:} value at the base point \(f(x)\).

\textbf{First order:} hyperplane tangent to the graph of \(f\):
\[
f_{\text{lin}}(x+h)=f(x)+\nabla f(x)^\top h.
\]

\textbf{Second order:} quadratic bowl approximation:
\[
f_{\text{quad}}(x+h) =
f(x)+\nabla f(x)^\top h + \tfrac12 h^\top H(x)h.
\]

\begin{itemize}
  \item If \(H(x)\succ 0\) → function locally convex ``bowl''.
  \item If \(H(x)\prec 0\) → concave ``cap''.
  \item Saddle behavior when Hessian is indefinite.
\end{itemize}

\textbf{Link to optimization:} classification of stationary points via Hessian eigenvalues.

\end{frame}

\begin{frame}{Remainder Term (Multivariate)}

\begin{definitionblock}{Lagrange-Type Remainder (informal)}
If \(f\in C^{2}\) and \(H\) is continuous near \(x\), then
\[
R_3(h) 
= \frac{1}{6}\sum_{i,j,k}
\partial_{ijk} f(x + \theta h)\, h_i h_j h_k,
\qquad \theta\in(0,1).
\]
\end{definitionblock}

\begin{itemize}
  \item Guarantees that the quadratic model is accurate for small \(h\).
  \item Useful in Newton method convergence analysis.
  \item In convex optimization: controls deviation from local quadratic model.
\end{itemize}

\end{frame}

% ======================================================
\subsection{Classical Inequalities \& Mean Value Results}
% ======================================================

\begin{frame}{Mean Value Theorem \& Rolle's Theorem}

\begin{definitionblock}{Rolle's Theorem}
If \(f\) is continuous on \([a,b]\), differentiable on \((a,b)\),
and \(f(a)=f(b)\), then \(\exists c\in(a,b)\) such that
\[
f'(c)=0.
\vspace{-3mm}
\]
\vspace{-3mm}
\end{definitionblock}

\begin{definitionblock}{Mean Value Theorem}
Under the same regularity assumptions,
\[
f'(c)=\frac{f(b)-f(a)}{b-a}
\qquad\text{for some }c\in(a,b).
\vspace{-3mm}
\]
\vspace{-3mm}
\end{definitionblock}
\vspace{-3mm}
\begin{itemize}
  \item Interpret \(f'(c)\) as ``average rate of change equals instantaneous rate.’’
  \item Foundation for inequalities, Lipschitz bounds, error estimates.
\end{itemize}

\end{frame}

\begin{frame}{Consequences of MVT}

\begin{itemize}
  \item \textbf{Bounded derivative implies Lipschitz continuity:}
    If \(|f'(x)| \le L\) then
    \[
    |f(x)-f(y)| \le L |x-y|.
    \]

  \item \textbf{Monotonicity:}
    \[
    f'(x)\ge 0\ \forall x \rightarrow f\text{ is increasing}.
    \]

  \item \textbf{Convexity:} 
    If and only if \(f'\) increasing if and only if \(f''\ge 0\) (when exists).

  \item \textbf{Error control via Taylor:}
    If \(|f^{(n+1)}|\le M\),
    \[
    |R_{n+1}(x)|\le 
    \frac{M}{(n+1)!}|x-x_0|^{n+1}.
    \]
\end{itemize}

Examples:
\begin{itemize}
  \item \(\log(1+x) \le x\) from concavity of \(\log\).
  \item \(\exp(x)\ge 1+x\) from convexity of \(\exp\).
\end{itemize}

\end{frame}

\begin{frame}{Classical Inequalities from Taylor/MVT}

\textbf{Bernoulli inequality:} \((1+x)^r \le 1+rx\) for \(r\in(0,1)\).

\textbf{Logarithm bounds:}
\[
x - \frac{x^2}{2} \le \log(1+x) \le x \quad (x>-1).
\]

\textbf{Exponential bounds:}
\[
1 + x \le e^x \le 1 + x + \frac{x^2}{2} e^{|x|}.
\]

\textbf{Sine inequality:}
\[
|\sin x| \le |x|.
\]

\textbf{All obtained from:}
\begin{itemize}
  \item MVT (monotonicity/convexity arguments),
  \item Taylor expansions with sign information on derivatives.
\end{itemize}

\end{frame}


% ==================================================
% SECTION 4 — Links to Optimization: Convexity \& Hessian
% ==================================================
% ======================================================
\section{Optimization Links: Convexity \& the Hessian}
% ======================================================

\begin{frame}{Outline}
  \tableofcontents[currentsection,hideothersubsections]
\end{frame}

% ------------------------------------------------------
\subsection{First/Second Order in Optimization}
% ------------------------------------------------------

\begin{frame}{Stationarity \& Quadratic Approximation}

\begin{definitionblock}{First-Order Optimality Condition}
For an unconstrained differentiable function \(f:\mathbb{R}^n\to\mathbb{R}\),
\[
x^\star \text{ is a local extremum } \Rightarrow \nabla f(x^\star)=0.
\]
\end{definitionblock}

\begin{itemize}
  \item The gradient \(\nabla f(x)\) gives the direction of steepest ascent.
  \item At a stationary point, all first-order directional derivatives vanish.
  \item Near \(x^\star\), approximate \(f\) by its \textbf{local quadratic model:}
  \[
  m(h)=f(x^\star)+\tfrac{1}{2}h^\top\nabla^2 f(x^\star)h.
  \]
\end{itemize}

\end{frame}

% ------------------------------------------------------

\begin{frame}{Second-Order Classification (C\(^2\) Functions)}

\begin{definitionblock}{Second-Order Test (Unconstrained Case)}
If \(f\in C^2\) and \(\nabla f(x^\star)=0\), then:
\[
\begin{cases}
\nabla^2 f(x^\star) \succ 0 & \Rightarrow \text{local minimum},\\[0.3em]
\nabla^2 f(x^\star) \prec 0 & \Rightarrow \text{local maximum},\\[0.3em]
\nabla^2 f(x^\star) \text{ indefinite} & \Rightarrow \text{saddle point.}
\end{cases}
\vspace{-6mm}
\]
\end{definitionblock}

\begin{itemize}
  \item Eigenvalues of the Hessian determine the curvature:
  \[
  f(x+h)\approx f(x^\star)+\tfrac{1}{2}\sum_i \lambda_i (u_i^\top h)^2.
  \]
  \item Positive curvature in all directions → local minimum.
  \item Mixed curvatures → saddle.
\end{itemize}

\end{frame}

% ------------------------------------------------------
\subsection{Convexity via the Hessian}
% ------------------------------------------------------

\begin{frame}{Convexity Test (C\(^2\))}

\begin{definitionblock}{Convexity \& Hessian Criterion}
If \(f\in C^2\) on a convex domain and 
\[
\nabla^2 f(x)\succeq 0,\ \forall x,
\]
then \(f\) is convex.  
If \(\nabla^2 f(x)\succ 0\) (positive definite), then \(f\) is \textbf{strictly convex}.
\end{definitionblock}

\begin{itemize}
  \item \textbf{Geometric meaning:} tangent planes lie below the graph of \(f\).
  \item \textbf{Implication:} any stationary point is the \emph{global} minimizer.
  \item \textbf{Stronger:} if \(\nabla^2 f(x)\succeq \mu I\) (\(\mu>0\)) → \emph{strong convexity}.
\end{itemize}

\[
f(y)\ge f(x)+\nabla f(x)^\top(y-x)+\frac{\mu}{2}\|y-x\|^2.
\]

\end{frame}

% ------------------------------------------------------
\begin{frame}{Convexity \& Global Optimality}

\textbf{First-order convexity test:}
\[
f(y)\ge f(x)+\nabla f(x)^\top(y-x)
\quad \forall x,y.
\]
\textbf{Equivalent (for \(C^2\) functions):}
\[
\nabla^2 f(x)\succeq 0 \ \forall x.
\]

\begin{itemize}
  \item For convex \(f\): any point with \(\nabla f(x^\star)=0\) is a \textbf{global minimizer}.
  \item If \(f\) is strictly convex: minimizer is \textbf{unique}.
  \item Strong convexity provides quadratic lower bounds → crucial for convergence rates in optimization algorithms (e.g., gradient descent).
\end{itemize}

\end{frame}

% ------------------------------------------------------
\begin{frame}{Examples of Convex/Non-Convex Functions}

\textbf{Convex:}
\[
\begin{aligned}
&f(x)=x^2 &&\Rightarrow \nabla^2 f(x)=2>0,\\
&f(x)=e^x &&\Rightarrow f''(x)=e^x>0,\\
&f(x)=\|Ax-b\|_2^2 &&\Rightarrow \nabla^2 f(x)=2A^\top A\succeq 0.
\end{aligned}
\]

\textbf{Non-convex:}
\[
f(x,y)=x^2-y^2
\Rightarrow 
\nabla^2 f(x,y)=
\begin{pmatrix}
2 & 0\\[2pt]
0 & -2
\end{pmatrix}
\text{ (indefinite)}.
\]

\textbf{Implications:}
- Convex imply all local minima are global.  
- Non-convex imply saddle points or local minima possible.

\end{frame}

% ------------------------------------------------------
\begin{frame}{Summary: Role of the Hessian in Optimization}

\begin{itemize}
  \item The \textbf{gradient} identifies candidate stationary points.
  \item The \textbf{Hessian} classifies curvature and determines convexity.
  \item \textbf{Convexity} ensures global optimality and algorithmic stability.
  \item \textbf{Strong convexity} → faster convergence of gradient-based methods.
  \item Used extensively in:
    \begin{itemize}
      \item Portfolio optimization (quadratic risk functions),
      \item Calibration of least squares,
      \item Sensitivity analysis (Taylor–Hessian local model).
    \end{itemize}
\end{itemize}

\[
\boxed{
\nabla f(x^\star)=0,\quad
\nabla^2 f(x^\star)\succ 0
\Rightarrow
x^\star=\text{local (global if convex) minimum.}
}
\]

\end{frame}



\end{document}